\documentclass[11pt]{article}
\usepackage[margin=1in]{geometry}
\usepackage[utf8]{inputenc} % Added for safety
\usepackage[T1]{fontenc}    % Good practice for font encoding
\usepackage{amsmath,amssymb,amsthm}
\usepackage{graphicx}
\usepackage{booktabs}
\usepackage{caption}
\usepackage{listings}
\usepackage{enumitem}
\usepackage{hyperref}

\title{Image-based Scagnostic Pipeline (Subpixel / Continuous Geometry)}
\author{Generated pipeline}
\date{\today}

\begin{document}
\maketitle

\begin{abstract}
This document gives a concise, self-contained description of an image-only scagnostics pipeline designed to reproduce the geometric intent of traditional graph-based scagnostics (Wilkinson et al.) while remaining purely raster-based. It contains algorithmic steps, continuous formulas and their discrete approximations, parameter suggestions, and a mapping from traditional graph components to image-based replacements.
\end{abstract}

\section{Overview}
Assumption: the input image \(I(x,y)\) is a sampled representation of a continuous intensity/density \(f(x,y)\). The pipeline reconstructs continuous geometric primitives (subpixel contours, medial axis with radius) and computes Euclidean measures (area, perimeter, arc-length) that mimic graph-based scagnostic quantities.

High-level steps:
\begin{enumerate}[leftmargin=*,noitemsep]
  \item Input \& anti-alias (smooth)
  \item Percentile / multi-threshold segmentation
  \item Subpixel contour extraction (marching-squares)
  \item Optional contour smoothing (spline)
  \item Continuous geometry: area, perimeter, convex hull
  \item Distance transform \(\rightarrow\) medial axis (continuous skeleton)
  \item Weighted skeleton integrals and branch stats
  \item Orientation field via structure tensor
  \item Multi-scale aggregation \& normalization
\end{enumerate}

\section{Notation}
\begin{description}
  \item[\(I\)] input raster image (floating values)
  \item[\(\sigma\)] Gaussian smoothing standard deviation (pixels)
  \item[\(T\)] a chosen threshold (often a percentile of \(I\))
  \item[\(M\)] binary mask \(M(x,y) = \mathbf{1}_{I(x,y) \ge T}\)
  \item[\(C\)] subpixel contour: ordered list of floating-point vertices \(\{(x_i,y_i)\}_{i=1}^{N}\)
  \item[\(P\)] perimeter (continuous approximation)
  \item[\(A\)] area (continuous, polygon area)
  \item[\(d(x,y)\)] Euclidean distance transform (distance to background)
  \item[\(\mathcal{S}\)] medial axis / skeleton represented as continuous centerlines
  \item[\(s\)] arc-length parameter along \(\mathcal{S}\)
  \item[\(\kappa\)] curvature where needed
\end{description}

\section{Step-by-step pipeline and formulas}

\subsection{1. Anti-alias / smoothing}
Apply light Gaussian blur to stabilize interpolation:
\[
I_s = G_{\sigma} * I,\qquad \sigma \in [0.5,2.0]\ \text{pixels (typical)}
\]
This reduces quantization and makes marching-squares interpolation stable.

\subsection{2. Percentile / multi-threshold segmentation}
Choose one or several thresholds \(T_k\) using percentiles of \(I_s\) (e.g., 90\%, 95\%, 99\%) to create masks:
\[
M_k(x,y) = \mathbf{1}_{I_s(x,y) \ge T_k}.
\]
Multi-thresholds emulate the family of \(\alpha\)-shapes.

\subsection{3. Subpixel contour extraction (marching squares)}
For each mask \(M_k\), extract contours using marching squares / linear interpolation to produce a polyline \(C = \{(x_i,y_i)\}\) with floating coordinates. The contour vertices are subpixel coordinates.

\subsection{4. Contour smoothing (optional)}
If required, fit a smoothing spline \(\mathbf{c}(t)\), \(t\in[0,1]\), to the polyline points to remove high-frequency noise (e.g., cubic B-spline or Catmull--Rom).
Use smoothing parameter chosen to preserve small features; avoid oversmoothing.

\subsection{5. Continuous area and perimeter}
Given a closed polyline \(C=\{(x_i,y_i)\}_{i=1}^N\) (ordered, closed):
\paragraph{Area (shoelace formula)}
\[
A = \frac{1}{2}\left| \sum_{i=1}^N (x_i y_{i+1} - x_{i+1} y_i)\right|,\qquad (x_{N+1},y_{N+1})=(x_1,y_1).
\]
\paragraph{Perimeter (polyline)}
\[
P = \sum_{i=1}^{N} \sqrt{(x_{i+1}-x_i)^2 + (y_{i+1}-y_i)^2}.
\]
If using a spline \(\mathbf{c}(t)\), compute
\[
P = \int_{0}^{1} \|\mathbf{c}'(t)\|\,dt
\]
and approximate numerically.

\paragraph{Isoperimetric quotient (for ``Skinny'')}
\[
\mathrm{IQ} = \frac{4\pi A}{P^2},\qquad \text{Skinny}\approx 1-\mathrm{IQ}\ (\text{or normalized}).
\]

\subsection{6. Convex hull (subpixel)}
Compute convex hull of contour vertices \(\mathrm{hull}(C)\). Let \(A_{\mathrm{hull}}\) be its area. Convex measure:
\[
\text{Convex} = \frac{A}{A_{\mathrm{hull}}}.
\]

\subsection{7. Distance transform and medial axis (continuous skeleton)}
Compute Euclidean distance transform \(d(x,y)\) on the binary mask \(M\):
\[
d(x,y) = \min_{(u,v)\in \text{background}} \sqrt{(x-u)^2 + (y-v)^2}.
\]
Extract medial axis \(\mathcal{S}\) as the ridge of \(d\) (maxima of distance along perpendicular directions). Represent \(\mathcal{S}\) as ordered centerline curves with associated local radius \(d(s)\).

\paragraph{Discrete approximation for skeleton arc-length}
If skeleton points along \(\mathcal{S}\) are \(\{(x_i,y_i)\}\), arc-length:
\[
L \approx \sum_{i} \sqrt{(x_{i+1}-x_i)^2 + (y_{i+1}-y_i)^2}.
\]

\paragraph{Weighted skeleton integral (MST-like)}
To mimic spacing-sensitive MST edge-lengths, use:
\[
L_w \approx \sum_{i} d(x_i,y_i)\cdot \Delta s_i
\]
where \(\Delta s_i\) is segment length. This weights the centerline by local width (analogous to edge-length \(\times\) spacing).

\subsection{8. Branch statistics and outlying}
Extract connected branches of \(\mathcal{S}\) and compute branch lengths \(\{\ell_j\}\) and radii \(\{r_j\}\).
Define outlying score as fraction of total mass (or area) in components/branches with small radius:
\[
\text{Outlying} \approx \frac{\sum_{j: r_j < r_{\mathrm{th}}} A_j}{A},
\]
or the fraction of branches with \(\ell_j\) much smaller/larger than median branch length (tunable threshold).

\subsection{9. Clumpy (density blobs)}
Compute watershed segmentation / connected components on smoothed density peaks to find blobs. Let \(B\) be number of blobs and \(a_i\) their areas. Example clumpiness measure:
\[
\text{Clumpy} = \frac{B\cdot \operatorname{Var}(\{a_i\})}{A^2} \quad \text{(normalized)}
\]

\subsection{10. Orientation field and striated / monotonic}
Compute image gradients \(I_x,I_y\) (e.g., Sobel) and the structure tensor at each pixel:
\[
J = \begin{pmatrix} I_x^2 & I_x I_y \\ I_x I_y & I_y^2 \end{pmatrix},
\]
optionally smoothed with a Gaussian window. Let eigenvalues be \(\lambda_1\ge\lambda_2\). Coherency (local orientation strength):
\[
\text{coh} = \frac{\lambda_1 - \lambda_2}{\lambda_1 + \lambda_2 + \epsilon}.
\]
Collect orientation angles \(\theta(x,y)=\tfrac{1}{2}\arctan2(2I_xI_y, I_x^2 - I_y^2)\) in the mask and compute circular variance; striated score \(\approx 1-\) circular variance. For monotonicity, sample points along the principal skeleton path and compute Spearman correlation \(\rho\) between x and y ranks:
\[
\text{Monotonic} \approx |\rho|.
\]

\subsection{11. Multi-scale aggregation and normalization}
Repeat key metrics across \(k\) thresholds \(\{T_k\}\) or radii \(r\), then aggregate (mean, median, slope across scales). Finally, robust-scale each metric (subtract median, divide by IQR) or apply rank-preserving mapping (isotonic regression) for downstream tasks.

\section{Discrete approximations and numerical notes}
\begin{itemize}
  \item \textbf{Marching squares}: find level contours at 0.5 on smoothed mask \(M\) or at chosen intensity thresholds on \(I_s\). Resulting vertex coordinates are floating.
  \item \textbf{Spline arc-length}: discretize spline at fine parameter resolution (e.g., 100--500 samples) and sum segment lengths.
  \item \textbf{Distance transform}: use Euclidean distance transform (EDT) and ridge extraction (e.g., non-maximum suppression along gradient direction).
  \item \textbf{Skeleton pruning}: remove spur branches shorter than a fraction of image diagonal (e.g., \(0.5\%\)--\(2\%\) diag).
  \item \textbf{Stability}: use multi-scale smoothing and percentile thresholds to avoid brittle decisions.
\end{itemize}

\section{Mapping to original scagnostic objects}
\begin{table}[h]
\centering
\begin{tabular}{@{}p{0.32\linewidth}p{0.32\linewidth}p{0.32\linewidth}@{}}
\toprule
\textbf{Original graph object / purpose} & \textbf{Image-only replacement} & \textbf{Why (intent match)} \\
\midrule
Delaunay / $\alpha$-shape ($\alpha$-hull) & Subpixel contour(s) at multi-thresholds & Contours at varying levels emulate $\alpha$-shape family; vertices allow continuous hull area. \\
MST (edge lengths) & Medial axis (distance-aware skeleton) & Skeleton centerlines + local radius approximate centerline connectivity and spacing. \\
Edge-length statistics / Outliers & Branch length \& radius distribution & Long MST edges appear as long thin branches or large empty basins in distance transform. \\
Convex hull / compactness & Convex hull of contour vertices & Polygon area and hull area computed continuously match convexity intent. \\
Perimeter-based skinny & Spline perimeter \& area \(\rightarrow\) IQ & Continuous perimeter avoids staircase bias; IQ maps to skinny measure. \\
Clump detection & Watershed / blob segmentation on density & Number and variance of blobs reproduce clumpy notion. \\
Striated / orientation & Structure tensor coherence + long narrow branches & Orientation coherence corresponds to aligned edges/striations. \\
Monotonic / trend & Spearman on principal skeleton path samples & Principal path ranking reproduces monotonic ordering. \\
\bottomrule
\end{tabular}
\caption{Practical mapping between graph-based concepts and image-based replacements.}
\end{table}

\section{Algorithm (pseudocode)}
\begin{lstlisting}[language=Python, basicstyle=\ttfamily\small]
# Input: image I (float), scale_params, thresholds T_k
I_s = gaussian_filter(I, sigma=sigma)
for T in thresholds:
    M = (I_s >= percentile(I_s, T))
    contours = marching_squares(M)  # subpixel polylines
    main_contour = select_largest_contour(contours)
    if smooth_contour:
        c = fit_spline(main_contour)
    else:
        c = main_contour
    A = polygon_area(c)
    P = polyline_length(c) or spline_arc_length(c)
    hull = convex_hull(c)
    Convex = A / area(hull)

    d = euclidean_distance_transform(M)
    S = extract_medial_axis(d)    # ridge extraction
    S = prune_short_branches(S, min_len=0.005*diag)
    L = arc_length(S)
    Lw = sum( d(p) * ds for p along S )
    branch_stats = compute_branch_lengths_and_radii(S, d)

    blobs = watershed_on_density(I_s)
    clumpy = measure_blob_count_and_variance(blobs)

    J = structure_tensor(I_s)
    coh_map = coherency(J)
    striated = 1 - circular_variance(orientation_angles, mask=M)
    monotonic = abs(spearmanr(sample_points_along_main_path(S)))
# Aggregate metrics across thresholds, normalize, output vector
\end{lstlisting}

\section{Parameter recommendations (defaults to try)}
\begin{itemize}
  \item Raster resolution: $\ge$512 px on longest side if possible.
  \item Gaussian blur $\sigma$: 0.5--2.0 px (start with 1.0).
  \item Thresholds (percentiles): \{90, 95, 99\} or \{80, 90, 95\} depending on contrast.
  \item Skeleton pruning length: 0.5\%--2.0\% of image diagonal.
  \item Contour spline samples for arc-length: 200 samples per closed contour.
  \item Structure tensor smoothing window: 3--9 px depending on scale.
  \item Normalization: robust z-score (subtract median, divide by IQR).
\end{itemize}

\section{Validation checklist (recommended)}
\begin{enumerate}[leftmargin=*,noitemsep]
  \item Compare rank correlations (Spearman) between image metrics and graph metrics on synthetic canonical families (e.g., line, cluster, sparse, uniform).
  \item Visual inspection: overlay subpixel contours \& skeleton on image.
  \item Stability: vary $\sigma$, thresholds, and pruning length; measure metric variance.
  \item Clustering sanity: run clustering on image-scagnostics vs graph-scagnostics and compute ARI.
\end{enumerate}

\section{Notes and caveats}
\begin{itemize}
  \item This pipeline is intended for images that represent densities or rasterized scatterplots (KDEs, heatmaps). It is \emph{not} appropriate for symbolic categorical rasters, low-resolution pixel art, or images where pixel adjacency is the only meaningful relation.
  \item Even with continuous approximations, image scagnostics approximate the graph-based values; they are typically closer in rank/ordering than in raw numeric equality. Use rank-preserving mappings (isotonic regression) if numeric equivalence is required.
\end{itemize}

\end{document}